\documentclass[../DoAn.tex]{subfiles}
\begin{document}

\begin{center}
    \Large{\textbf{TÓM TẮT NỘI DUNG ĐỒ ÁN}}\\
\end{center}
\vspace{1cm}

Máy chấm công bằng khuôn mặt phải chặn được ảnh giả và không nên gửi ảnh khuôn mặt ra khỏi máy, trong khi ba mô hình phát
hiện, chống giả mạo và nhận diện sau lượng tử hóa vẫn nặng hơn bốn lần SRAM của một vi điều khiển. Đồ án xây dựng một hệ
thống chấm công gồm máy chấm công trên ESP32-S3, máy chủ quản trị nhân sự và cổng nhân viên trên điện thoại.

Máy chấm công chạy YuNet, MiniFASNetV1SE và MobileFaceNet-ECA trên thư viện ESP-DL. Mỗi nhánh đi qua một thang lượng tử hóa
bốn bậc, và cấu hình được chọn bằng thước của chính máy. Các phép toán tự viết bằng lệnh vector của chip đưa một lượt ba mô
hình từ 444,8~ms xuống 333,8~ms mà đầu ra không đổi một bit. Bản đang chạy mất 353,6~ms lúc rảnh và 406,2~ms khi có xem
trước, nhanh hơn TFLite Micro 10,51 lần. Lượt chấm được ghi xuống flash trước rồi mới gửi qua MQTT QoS 1, và máy chủ bỏ
bản trùng bằng ràng buộc unique, nên mất mạng không làm mất lượt chấm. Máy chủ NestJS và PostgreSQL đi từ lượt chấm tới
ngày công, nghỉ phép, thuế thu nhập cá nhân, bảo hiểm và phiếu lương, kèm dashboard web, PWA và ứng dụng Android.

Trên 87 khung chụp bằng camera OV5640 của máy, nhánh chống giả mạo chặn 25/25 khung giả và giữ 59/62 mặt thật. Bộ phát hiện
đạt AP 0,9544 trên mặt từ 38~px. Backend qua 971/971 ca kiểm thử e2e. Phần chưa đạt được ghi cùng số đo: nhận diện sụt
3,7\% sau lượng tử hóa, chưa có khung ảnh in để đo và chưa đo FAR, FRR với nhiều người. Trên VPS 1 vCPU với dữ liệu 5.000
nhân viên, máy chủ chịu được 10 phiên đồng thời nhưng sập ở 20 phiên, và lượt tính lương cho 4.776 người chưa chạy xong.

\vspace{1cm}
\begin{flushright}
Sinh viên thực hiện\\
\begin{tabular}{@{}c@{}}
\textit{Đỗ Văn Việt}
\end{tabular}
\end{flushright}

\end{document}
